% -*- TeX-master: "../main.tex" -*-
\chapter{Podsumowanie}
Zaprojektowany algorytm genetyczny pozwolił na efektywne rozwiązanie problemu optymalizacji sieci komunikacyjnej dla miasta Sławków. Dzięki zastosowaniu wielokryterialnej funkcji fitness oraz mechanizmów zabezpieczających przed utknięciem w lokalnych ekstremach, uzyskane wyniki wykazują wysoką zgodność z założonymi celami projektowymi. Implementacja w środowisku .NET z wykorzystaniem biblioteki Raylib umożliwiła nie tylko sprawne obliczenia, ale również czytelną wizualizację wypracowanego rozkładu jazdy. 

Mimo złożoności problemu optymalizacyjnego, zespół projektowy dostarczył rozwiązanie, które stanowi użyteczny model wspomagający planowanie transportu miejskiego. Wypracowane wagi oraz operatory ewolucyjne mogą służyć jako baza do dalszych prac nad udoskonalaniem systemu komunikacji publicznej w przyszłości.
